\documentclass[12pt,a4paper]{article}

\input{/home/tabaka/Documents/GitHub/Defs/defs.tex}

\setcounter{zadaniaRozdzial}{4}
%\pagestyle{empty}

\begin{document}
	
	\begin{center}
		\LARGE Funkcja liniowa
	\end{center}
	
	\ZbiorZadanieTwoXThree{Naszkicuj podane proste w układzie współrzędnych:}
	{$y=2x-3$}{$y=5$}
	{$y=\frac{2}{3}x+1$}{$y=-x-2$}
	{$y=-\frac{4}{5}x+3$}{$y=1\frac{1}{2}x-4$}
	
	\ZbiorZadanieTwoXThree{Przez które ćwiartki układu współrzędnych przechodzi podana prosta?}
	{$y=3-2x$}{$y=\frac{1}{5}x+2$}
	{$y=\sqrt{3}x+\sqrt{2}$}{$y=(3\sqrt{2}-2\sqrt{3})x+(7-5\sqrt{2})$}{a}{a}
	
	\ZbiorZadanieTwoXThree{Wyznacz równanie prostej przechodzącej przez punkty $A$ i $B$:}
	{$A=(-3,4)\qquad B=(1,0)$}{$A=(5,-1)\qquad B=(3,3)$}
	{$A=(5,6)\qquad B=(1,4)$}{$A=(2,3)\qquad B=(-1,-3)$}
	{$A=(4,-2)\qquad B=(-2,1)$}{$A=(4,-4)\qquad B=(-2,-4)$}
	
	\ZbiorZadanieTwoXThree{Określ monotoniczność podanej funkcji liniowej:}
	{$y=\frac{2}{3}x-4$}{$y=-2x+5$}
	{$y=-7+8x$}{$y=\frac{1-\sqrt{2}}{3}x-4$}{a}{a}
	
	\ZbiorZadanieTwoXThree{Wyznacz parametr $m$, dla którego prosta jest malejąca:}
	{$y=(m-3)x+7+m$}{$y=5x+m+mx$}
	{$y=\frac{m+8}{5}x+\frac{3m-1}{7}$}{$y=(8-5m)x+\sqrt{3}m$}{}{}
	
	\ZbiorZadanieTwoXThree{Wyznacz parametr $k$, dla którego podana prosta przechodzi przez punkt $P=(2,-3)$:}
	{$y=2kx-3$}{$y=(2k+1)x+3k+2$}
	{$y=(2k-1)x+4x-5$}{$y=\frac{5k+2}{6}x-2k+1$}
	
	\newpage \newpage
	
	\begin{center}
		\LARGE Zadania maturalne
	\end{center}
	

	
	\ZadanieMaturalneABCD{Dana jest funkcja liniowa określona wzorem
		$$y=(m^2-4)x+m-2.$$
		Funkcja ta nie ma miejsc zerowych, gdy}
	{$m=2$}
	{$m=-2$}
	{$m=0$}
	{$m=\sqrt{2}$}
	
	
	
\end{document}
